\chapter{Theory of Everything}
\label{chap:theory-of-everything}


\section*{The Ultimate Movie Theater}

Let us imagine an ultimate movie theater in which the film itself is encoded using a complex-valued spectral codec.

Unlike conventional cinema, the film is not projected onto a flat, two-dimensional rectangular screen. Its geometric representation is a three-dimensional curved manifold, with no predetermined topology.

Nor does this representation have a uniform pixel distribution. Its informational resolution varies spatially, concentrating representational resources in regions where greater detail is required.

Furthermore, the relationship between the spectral encoding and its geometric representation is not unique. The same spectral information can admit multiple geometric interpretations,
while the same geometric structure can be represented through multiple spectral encodings.

The question then becomes: which geometric representation does an observer actually experience?

The proposed answer is that observer-compatible representations with the greatest joint spectral-geometric compression are statistically favored.

In other words, an observer is most likely to find themselves embedded in a geometry that provides the most economical joint description of the spectral information and its geometric manifestation.


\subsection*{The mathematical formulation}


Let \(\psi\) denote the spectral representation and \(G\) a geometric representation.

We seek to minimize

\[
C\_{\mathrm{total}}(\psi,G) = C\_{\mathrm{spectral}}(\psi) + C\_{\mathrm{geometric}}(G)
\]

subject to the existence of an observer-compatible interpretation.

More generally, allowing for interactions between the two representations:

\[
C\_{\mathrm{total}}(\psi,G) = C\_{\psi}(\psi) + C_G(G) + C\_{\mathrm{interaction}}(\psi,G)
\]

The observer-conditioned probability distribution could then be expressed as

\[
P(G,\psi\mid O) = \frac{1}{Z_O} \exp\left[-\lambda C\_{\mathrm{total}}(\psi,G)\right]
\]

Here, the observer condition is essential: we are not simply looking for the mathematically simplest geometry, but the most economical joint representation compatible with an observer.


The proposal does not predict that the universe must have the simplest possible geometry.
It predicts that the geometry experienced by an observer is selected from a much larger space of possible representations according to joint compression and observer compatibility.

A geometrically complicated universe might actually be the most economical representation of a particular spectral structure.

Conversely, a geometrically simple universe might require an extremely complicated spectral description.

The optimum is determined by the combination, not either component independently.

The spacetime geometry is neither an arbitrary background nor necessarily a fundamental entity.
It is a representation that participates in the joint compression of physical information.




\section{Mutual Spectral-Geometric Compression}

General relativity couples energy-momentum to spacetime curvature.
In the present framework, a deeper relationship is proposed: spectral and geometric descriptions mutually constrain their respective compression.

Geometric complexity constrains which spectral representations provide economical descriptions, while spectral complexity constrains which geometric structures provide economical representations.

The resulting geometry is therefore not selected from a predetermined collection of possible spacetime structures.
Its shape and topology emerge together with the spectral representation that makes their joint description most economical.

This mutual optimization also provides a way to understand the emergence of classical geometry through equivalence classes.

Consider an hourglass. An astronomical number of microscopic sand-grain configurations produce approximately the same macroscopic conical pile.
Although a configuration resembling the Mona Lisa is theoretically conceivable, the number of microscopic arrangements compatible with such a
structure is negligible compared with those producing an ordinary conical pile.

Cone therefore should not wonder why its happens to have such a conical shape. The cone is therefore the largest not a particular microscopic arrangement, but a macroscopic equivalence class containing an enormous number of underlying configurations.

The same reasoning can be extended to spectral and geometric representations. Many distinct spectral configurations can correspond to essentially the same observable geometry, while multiple geometric structures can provide equivalent descriptions of the same spectral information.

An observer does not necessarily distinguish between these underlying alternatives. What appears as a definite physical reality can therefore correspond to an extensive equivalence class of compatible spectral-geometric configurations.

The proposed principle is that observer-compatible equivalence classes with the greatest joint statistical weight dominate the observed physical descriptions.

In this interpretation, quantum mechanics and general relativity are not independent fundamental descriptions. They are mutually constraining aspects of a joint spectral-geometric compression process.


\section{The Emergence of Resolution}

What determines the size, density and spatial distribution of the universe's pixels?

The universe does not have pixels arranged in a particular way because somebody designed the screen. 

Nothing determines them in advance. They are emergent properties of the underlying spectral-geometric configurations.

Consider the hourglass again. All possible arrangements of sand grains are conceivable, but the vast majority of configurations compatible with its macroscopic constraints produce approximately the same cone-shaped pile.
The cone is not an imposed structure or an instruction followed by individual grains. It is the macroscopic manifestation of an enormous statistical equivalence class.

Now imagine an observer whose existence depends on the macroscopic structure of the hourglass. If the cone-shaped configuration represents the overwhelmingly dominant observer-compatible class, the observer is most likely to find themselves in a cone-shaped universe.

The same principle applies to the universe.

There is no predefined pixel grid, uniform spatial resolution or prescribed geometric topology. All possible distributions are admitted, but they do not necessarily have equal statistical weight.
Some spectral-geometric configurations correspond to vastly larger observer-compatible equivalence classes than others.

The spatial distribution of matter, the resolution of spacetime and its geometric structure are therefore not independently prescribed.
They emerge together as properties of the statistically dominant observer-compatible configurations.

The universe does not need to construct a particular screen. The screen itself is an emergent property of the configurations that dominate the observer-conditioned description.


